\documentclass[11pt,a4paper]{article}
\usepackage[T1]{fontenc}
\usepackage[utf8]{inputenc}
\usepackage[english]{babel}
\usepackage{lmodern,microtype}
\usepackage[a4paper,margin=2.55cm]{geometry}
\usepackage{amsmath,amssymb,amsthm,mathtools}
\usepackage{booktabs,tabularx,array,xurl}
\usepackage[colorlinks=true,linkcolor=blue!45!black,urlcolor=blue!50!black,citecolor=blue!45!black]{hyperref}
\usepackage{xcolor}
\setlength{\emergencystretch}{3em}
\newtheorem{theorem}{Theorem}[section]
\newtheorem{proposition}[theorem]{Proposition}
\newtheorem{lemma}[theorem]{Lemma}
\theoremstyle{remark}
\newtheorem{remark}[theorem]{Remark}
\newcolumntype{Y}{>{\raggedright\arraybackslash}X}
\title{\textbf{Strong-Coupling Lattice Gauge Theory:\\
Corrections to Poincar\'e--Dobrushin Bounds and\\
Conditional Gap-Transport Criteria}\\[4pt]
\large Corrected interim edition, version 2.7}
\author{Lukas Geiger\\
\small Independent Researcher, Bernau, Germany\\
\small ORCID: \href{https://orcid.org/0009-0005-7296-1534}{0009-0005-7296-1534}}
\date{21 November 2026}
\begin{document}
\maketitle
\thispagestyle{empty}
\begin{center}
\fcolorbox{black!55}{black!4}{\parbox{0.88\linewidth}{\centering
\textbf{CORRECTED PREPRINT / REPLACEMENT INTERIM EDITION}\\
This version withdraws the earlier volume-uniform physical
transfer-gap claim \emph{as proved by the original argument}.
It does not establish the four-dimensional continuum Yang--Mills mass gap.
The earlier full manuscript remains identifiable through
\href{https://doi.org/10.5281/zenodo.20716608}{Zenodo v2.6}.
}}
\end{center}
\begin{abstract}
This corrected interim edition revisits the claim in the preceding
FST--YM preprint that a volume-independent physical transfer gap follows
directly from Poincar\'e--Dobrushin machinery for the Wilson lattice
measure. That deduction is not established: the historical influence
row sum counted six rather than up to eighteen neighboring link
incidences, the restricted-typical-set argument used an invalid
concentration and restriction step, and the gap of a heat-bath
generator was not quantitatively identified with the physical
time-transfer operator. We state a conservative Dobrushin uniqueness
criterion and prove an elementary positive but spatial-volume-dependent
finite-lattice transfer bound for a standard strictly positive Wilson
transfer kernel. We also give a correctly normed abstract Schur
comparison and an explicitly conditional summable-loss criterion.
Neither supplies the missing uniform Yang--Mills RG estimates, a
noncircular infrared anchor, or a nontrivial
Osterwalder--Schrader continuum construction. The old asserted global
closure is withdrawn, not repaired by synthetic numerical ledgers.
\end{abstract}
\noindent\textbf{Keywords:} Yang--Mills; lattice gauge theory; correction;
Dobrushin uniqueness; transfer operator; spectral gap; renormalisation.\\
\textbf{MSC 2020:} 81T13, 82B20, 60K35.
\clearpage

\section{Purpose, provenance, and explicit correction map}
The original published version 2.6 presented a stronger conclusion
than was warranted by its calculations. The present document is a
\emph{replacement scholarly correction}, not an additional claim to
solve the Clay problem. The existence of a nontrivial four-dimensional
pure Yang--Mills theory with a positive physical mass gap remains an
open problem in the sense of Jaffe and Witten~\cite{JaffeWitten}.
Historical formulae and attempted extensions are preserved in the
previous version for traceability; they are not incorporated here as
valid theorems.
\begin{center}\small
\begin{tabularx}{\textwidth}{@{}p{0.23\textwidth}YY@{}}
\toprule
Previous assertion & Correction in this edition & Present status\\
\midrule
Six-neighbor Dobrushin row sum & Up to $6\cdot3=18$ link--plaquette incidences; conservative $\alpha\le18\tanh\beta$ & Conditional criterion for uniqueness, not a transfer gap\\
Heat-bath PI $\Rightarrow$ physical transfer gap & Different generators; explicit comparison is required & Unproved in the historical route\\
Typical-set improvement $e^{-O(\sqrt\beta)}$ & Chebyshev and restriction steps were invalid & Withdrawn\\
Uniform lattice gap from historical theorem & Elementary transfer minorisation gives a bound depending on spatial volume & Only finite spatial volume proved below\\
RG product as proof of a continuum gap & Requires actual relative operator bounds and uniform physical units & Abstract conditional lemma only\\
Osterwalder--Schrader continuum completion & Existence, nontriviality, complete reconstruction and mass scale missing & Open\\
\bottomrule
\end{tabularx}
\end{center}
The status map concerns the \emph{claims of the previous manuscript}.
Separate internal November 2026 investigations of sharper Dobrushin and
time-transfer estimates are not silently promoted to this revised
paper; their assumptions and independent review must be tracked in
their own proof records.

\section{Wilson normalisation and a corrected influence count}
Fix a four-dimensional hypercubic Wilson lattice without local
edge identifications, with gauge group $SU(N)$, $N\ge2$, and
fundamental coupling $\beta>0$:
\[
d\mu_{\Lambda,\beta}(U)=
Z^{-1}\exp\left[
  \frac{\beta}{N}\sum_{p\in P(\Lambda)}
      \operatorname{ReTr}(U_p)
\right]\prod_{e\in E(\Lambda)}dU_e.
\]
Here $dU_e$ is normalized Haar measure.
The action and the representation convention are those of
Wilson~\cite{Wilson}.

For an interior link $e$, exactly six plaquettes contain $e$.
Each such plaquette contains three \emph{other} links.
For the Dobrushin matrix
\[
c_{ee'}=\sup_{\substack{\omega,\omega'\\
\omega_f=\omega'_f\ \forall f\ne e'}}
\|\mu_e(\cdot\mid\omega)-\mu_e(\cdot\mid\omega')\|_{\rm TV},
\qquad \|\mu-\nu\|_{\rm TV}=\tfrac12\!\int|d\mu-d\nu|,
\]
the relevant row sum is over other \emph{links}, not over
plaquettes.

\begin{proposition}[Conservative row-sum estimate]\label{prop:count}
For a change in one neighboring link incident through one
plaquette, $c_{ee'}\le\tanh\beta$.
Counting incidences with multiplicity and using the triangle
inequality yields
\[
\boxed{\ \alpha(\beta):=\sup_e\sum_{e'\ne e}c_{ee'}
\le18\tanh\beta\ }.
\]
Thus $\beta<\operatorname{artanh}(1/18)$ is a sufficient
strong-coupling \emph{Dobrushin uniqueness} window for this
particular conservative estimate, subject to the stated local
geometry. This proposition does not assert a physical transfer gap.
\end{proposition}
\begin{proof}
Changing $e'$ modifies the effective one-link log weight by
$D(U)=\frac{\beta}{N}\operatorname{ReTr}(U(V-V')^\dagger)$.
Unitarity gives $|D(U)|\le2\beta$ and hence
$\operatorname{osc}D\le4\beta$.
For normalized exponential tilts with log-likelihood oscillation
$b$, the extremal two-point calculation gives total variation at
most $\tanh(b/4)$. This yields $\tanh\beta$ for one incident
plaquette. A link has at most eighteen other-link incidences, so
summing with multiplicity proves the row bound.
The standard Dobrushin criterion applies if $\alpha<1$.
Small periodic lattices with identified links must instead
use their actual multiplicities.
\end{proof}
Wu~\cite{Wu} gives functional inequalities under a Dobrushin
condition. Such an inequality belongs to its specified
Gibbs/heat-bath dynamics and \emph{does not by itself} identify the
physical Wilson time-transfer gap.

\section{What can be proved at fixed spatial volume}
Reflection positivity and strict positivity of the Wilson
transfer operator are classical lattice results, with the
physical gauge-invariant setting described by L\"uscher~\cite{Luscher}.
The following elementary quantitative argument is included so
that the smaller claim is stated without extrapolation.

Let $X=SU(N)^{E_s}$ for a \emph{fixed finite} spatial
lattice, with normalized product Haar measure $dx$.
Write $P_g$ for the orthogonal gauge-invariant projection,
$P_s$ for the spatial plaquette set and $E_s$ for the spatial
link set. Let
\[
W(x)=\exp\left[\frac\beta N\sum_{p\in P_s}
\operatorname{ReTr}(U_p(x))\right],
\quad
K(x,y)=\prod_{e\in E_s}
\exp\left[\frac\beta N\operatorname{ReTr}
(U_e(y)U_e(x)^{-1})\right],
\]
and
\[
T=P_g M_{\sqrt W}K M_{\sqrt W}P_g
     \big|_{\operatorname{ran}P_g}.
\]
The finite Wilson transfer realization must be constructed
with the stated gauge projection and reflection-positive time
slicing; the kernel formulas fix the finite setting in which
the bound is evaluated.

\begin{proposition}[Finite-spatial-volume transfer gap]\label{prop:finite}
For $\beta>0$, let $a_t>0$ be the time spacing. Suppose $T$
above is the physical strictly positive compact Wilson
transfer operator. Let $\rho$ be its Perron eigenvalue and
$H_{a_t}=-a_t^{-1}\log(T/\rho)$ on the vacuum-orthogonal
sector. Writing $D=|E_s|+|P_s|$, its positive
finite-volume gap satisfies
\[
\boxed{\ \Delta_{X,a_t}\ \ge
-\frac{1}{a_t}\log\!\left(1-e^{-4\beta D}\right)>0.\ }
\]
The lower bound depends explicitly on $D$ and is \emph{not}
uniform in growing spatial volume, coupling or continuum spacing.
\end{proposition}
\begin{proof}
Since $|\operatorname{ReTr}(U)|/N\le1$,
the continuous gauge-averaged kernel is bounded as
$0<m:=e^{-\beta D}\le T(x,y)\le M:=e^{\beta D}$.
Strict positivity and compactness yield a simple top
eigenvalue $\rho>0$ and normalized positive eigenfunction
$h$, with $\int_Xh^2dx=1$. Put $H=\int_Xh\,dx$ and define
the reversible Doob Markov kernel
\[
Q(x,dy)=\frac{T(x,y)h(y)}{\rho h(x)}\,dy,
\qquad \pi(dy)=h(y)^2\,dy,\qquad
\nu(dy)=h(y)\,dy/H.
\]
The eigenfunction equation gives
$mH\le\rho h(x)\le MH$. Therefore
$Q(x,\cdot)\ge(m/M)\nu(\cdot)$ and
$\nu\ge(m/M)\pi$, using $h_{\max}H\le M/m$.
Consequently $Q\ge(m/M)^2\pi$.
On mean-zero $L^2(\pi)$, the Markov contraction is at
most $1-(m/M)^2$; the strict operator positivity gives the
corresponding bound for the nonvacuum spectrum of $T/\rho$.
Taking the negative logarithm yields the display.
\end{proof}
At fixed volume this is a finite-dimensional-geometry
spectral fact, not a solution of the continuum mass-gap
problem. For $\beta\downarrow0$ the transfer dynamics
degenerates to the trivial Haar resampling limit.

\section{Why the typical-set proof must be withdrawn}
The original deduction combined a putative restricted
Poincar\'e constant with the probability of the discarded
set. It did not pay for the variance within disconnected
pieces or the exchange between the pieces. In general a
positive-measure restriction of a connected manifold need
not satisfy any gradient-form Poincar\'e inequality:
take two disjoint open subdomains and a function constant
with different values on each; its restricted gradient
energy vanishes while its variance is positive.
An estimate on the full measure therefore cannot be
restricted to an arbitrary typical set by multiplying
its constant by a power of its mass.

There is also an independent arithmetic error.
From
\[
\mu(|W-\mathbb EW|>\varepsilon)
\le\frac{C\beta^2}{\varepsilon^2},
\qquad \varepsilon=c\sqrt\beta,
\]
one obtains $C\beta/c^2$, \emph{not} $C/c^2$.
It does not give a uniformly small exceptional mass as
$\beta\to\infty$. Thus the previously stated improvement
from $e^{-O(\beta)}$ to $e^{-O(\sqrt\beta)}$ is withdrawn.
A valid repair would require geometry of the restricted
region, internal coercivity on \emph{both} sides, and a
separate lower bound on inter-region flow for the
\emph{same} physical Markov or transfer operator.

\section{Transfer, physical units, and the continuum}
A heat-bath Dirichlet form describes a resampling
dynamics; a reflection-positive time-transfer operator
describes the physical Euclidean evolution. Their gaps
are not interchangeable without an independently
established comparison of the respective quadratic
forms, ground states, and time normalisations.

If a positive $T_a$ has top eigenvalue $\rho_a$, the
physical generator is
\[
H_a=-a^{-1}\log(T_a/\rho_a).
\]
Its physical mass gap, if any, must remain bounded below
by a positive number along the actual joint
spatial-volume/UV-continuum trajectory.
A bound $\kappa(2a)\ge c\kappa(a)$
for some constant $c>0$ is not sufficient:
$\kappa(a)=a^2$ satisfies it with $c=4$ while
$\kappa(a)/a\to0$.

The Osterwalder--Schrader reconstruction programme
requires actual convergence of correctly normalised,
reflection-positive Schwinger functions and their
associated axioms~\cite{OS1,OS2}. Assuming uniform
exponential clustering in \emph{physical} distances
already assumes much of the required spectral input;
it is not an independent mass-gap construction.

\section{A correct abstract comparison and its limits}
The next two elementary lemmas provide possible
bookkeeping tools, \emph{not} verified Wilson RG estimates.

\begin{lemma}[Norm-corrected Schur comparison]\label{lem:schur}
Let $H\ge0$ be a finite-dimensional self-adjoint block
operator on a vacuum-orthogonal space,
\[
H=\begin{pmatrix}A&B\\B^*&C\end{pmatrix},
\quad C\ge MI>0,\qquad
S=A-BC^{-1}B^*,\quad G=I+BC^{-2}B^*.
\]
If $G^{-1/2}SG^{-1/2}\ge mI$ for some $m>0$,
then
\[
\boxed{\quad H\ge\frac{mM}{m+M}\,I.\quad}
\]
An unbounded-operator extension requires explicit
form-domain hypotheses.
\end{lemma}
\begin{proof}
For $u=(x,-C^{-1}B^*x)$ and
$v=(0,y+C^{-1}B^*x)$, the vector $(x,y)=u+v$
has energy
$\langle x,Sx\rangle+
\|C^{1/2}(y+C^{-1}B^*x)\|^2
\ge m\|u\|^2+M\|v\|^2$.
Weighted Cauchy--Schwarz bounds
$\|u+v\|^2\le(1/m+1/M)
(m\|u\|^2+M\|v\|^2)$.
\end{proof}

\begin{lemma}[Conditional summable-loss budget]\label{lem:budget}
Suppose each refinement step is comparably
normalised in \emph{physical energy units};
suppose $m_0>0$ is a volume- and cutoff-uniform
infrared anchor and a true effective-operator comparison
has relative log loss $r_j\ge0$, while the eliminated
sector has energy at least $M_j>0$. If
\[
\sum_{j\ge0}r_j<\infty,\qquad
\sum_{j\ge0}M_j^{-1}<\infty
\]
with uniform bounds across all finite-volume/cutoff
approximants, then the resulting gap certificates obey
\[
\boxed{\quad
m_n\ge
\frac{\exp(-\sum_{j<n}r_j)}
     {m_0^{-1}+\sum_{j<n}M_j^{-1}}>0.
\quad}
\]
\end{lemma}
\begin{proof}
Lemma~\ref{lem:schur} at step $j$ gives
$m_{j+1}^{-1}\le e^{r_j}m_j^{-1}+M_j^{-1}$.
Iterate and bound every later exponential multiplier
by $\exp(\sum_{k<n}r_k)$.
\end{proof}
If a quantitative running coupling $t_j=g_j^2$
satisfies $t_j-t_{j+1}\ge b_-t_j^2$
for $b_->0$, then
$\sum_{j\ge J}g_j^4\le g_J^2/b_-$
by telescoping. This only shows that
\emph{genuinely established} relative errors of
order $g_j^4$ would be summable. Neither
these relative operator errors nor an independent
uniform infrared anchor have been proved here
for four-dimensional Yang--Mills.
In the canonical strong-coupling Hamiltonian,
the expansion parameter $J/\kappa\sim g^{-4}$
instead diverges as $g\to0$; finite-order
strong-coupling calculations cannot simply be
carried into the asymptotically free ultraviolet.

\section{Corrected status and remaining proof obligations}
\begin{center}\small
\begin{tabularx}{\textwidth}{@{}p{0.32\textwidth}Y@{}}
\toprule
Status & Finding\\
\midrule
Finite spatial volume & Explicit strictly positive transfer bound, Proposition~\ref{prop:finite}; it worsens with volume.\\
Conservative small-$\beta$ control & Corrected 18-incidence Dobrushin criterion, Proposition~\ref{prop:count}.\\
Original volume-independent physical transfer-gap theorem & \textbf{Withdrawn as proved by the historical argument.}\\
Typical-set $e^{-O(\sqrt\beta)}$ improvement & \textbf{Withdrawn}; missing geometry and incorrect scaling.\\
Schur and summable-loss lemmas & Correct only under stated spectral and uniformity assumptions.\\
All-volume 4D physical gap, UV renormalisation and OS continuum & \textbf{Open}; not supplied by this revision.\\
\bottomrule
\end{tabularx}
\end{center}
The next scientifically meaningful work is a
noncircular, reflection-positive Yang--Mills
RG comparison with an independently proven
infrared gap at fixed physical units, uniform
coercivity of discarded sectors, summable
\emph{relative} form defects and an OS-valid
nontrivial continuum reconstruction. No
finite synthetic diagnostic can substitute
for these inputs.

\section{Outlook: prospective expanded Version 2.8}
Version~2.8 is planned as a \emph{fuller, independently re-examined}
development of the research programme, rather than a restoration of
the withdrawn theorem. Its intended subjects include exact
finite-lattice and two-plaquette calculations, the local
Schrieffer--Wolff and linked-cluster approaches, and a carefully
specified norm-corrected Schur/RG transport programme. Every
model-specific assertion will be separated from an independently
established \emph{physical} time-transfer comparison, uniform
infrared control and Osterwalder--Schrader continuum reconstruction.
Results will be included only after individual mathematical and
source checks. No publication date, nonperturbative continuum
construction or four-dimensional mass-gap proof is promised.

\section*{Version history and authorship disclosure}
\addcontentsline{toc}{section}{Version history and authorship disclosure}
\textbf{Version 2.7 (21 November 2026):}
replacement interim correction of the published
v2.6 text. Retracted unsupported
volume-uniform transfer and typical-set
claims; made the Dobrushin incidence count
explicit; retained a finite-volume bound
and conditional RG lemmas. No statement
of a proved four-dimensional continuum
Yang--Mills mass gap is made.
The earlier archival version remains
permanently accessible at
\href{https://doi.org/10.5281/zenodo.20716608}
{doi:10.5281/zenodo.20716608}.
\par\medskip
\textbf{AI-use disclosure.}
AI-assisted tools were used for
exploratory mathematics, drafting,
and computational checking.
The author is responsible for the
scientific content and publication
decisions. Tool-produced diagnostics
are not independent peer review.

\begin{thebibliography}{9}\small
\bibitem{JaffeWitten}
A.~Jaffe and E.~Witten,
\emph{Quantum Yang--Mills Theory},
in \emph{The Millennium Prize Problems},
Clay Mathematics Institute, 2000,
\url{https://www.claymath.org/wp-content/uploads/2022/06/yangmills.pdf}.
\bibitem{Wilson}
K.~G.~Wilson,
\emph{Confinement of quarks},
Phys.\ Rev.\ D \textbf{10} (1974), 2445--2459,
\href{https://doi.org/10.1103/PhysRevD.10.2445}{doi:10.1103/PhysRevD.10.2445}.
\bibitem{Luscher}
M.~L\"uscher,
\emph{Construction of a selfadjoint, strictly positive transfer
matrix for Euclidean lattice gauge theories},
Commun.\ Math.\ Phys.\ \textbf{54} (1977), 283--292,
\href{https://doi.org/10.1007/BF01614090}{doi:10.1007/BF01614090}.
\bibitem{Wu}
L.~Wu,
\emph{Poincar\'e and transportation inequalities for Gibbs
measures under the Dobrushin uniqueness condition},
Ann.\ Probab.\ \textbf{34} (2006), 1960--1989,
\href{https://doi.org/10.1214/009117906000000368}
{doi:10.1214/009117906000000368}.
\bibitem{OS1}
K.~Osterwalder and R.~Schrader,
\emph{Axioms for Euclidean Green's functions},
Commun.\ Math.\ Phys.\ \textbf{31} (1973), 83--112,
\href{https://doi.org/10.1007/BF01645738}{doi:10.1007/BF01645738}.
\bibitem{OS2}
K.~Osterwalder and R.~Schrader,
\emph{Axioms for Euclidean Green's functions II},
Commun.\ Math.\ Phys.\ \textbf{42} (1975), 281--305,
\href{https://doi.org/10.1007/BF01608978}{doi:10.1007/BF01608978}.
\end{thebibliography}
\end{document}
